O Fluxo de Potência Ótimo (FPO) é uma das ferramentas centrais para o planejamento e a operação de Sistemas Elétricos de Potência (SEP). No contexto brasileiro, programas consolidados de mercado, como o ANAREDE do CEPEL, incorporam ações de controle específicas (limites de potência reativa de geradores, ou QLIM; limites de magnitude de tensão em barras, ou VLIM; controle automático de susceptância de bancos chaveados, ou CSCA; controle automático de \textit{tap} de transformadores, ou CTAP; e, em emergência, esquemas hierárquicos de corte/alívio de carga, rotulados neste trabalho como DERA em referência ao nome do \textit{script} correspondente) que não estão presentes de forma nativa nas formulações acadêmicas clássicas implementadas em pacotes \textit{open-source}. Este trabalho apresenta a construção, em linguagem Julia, de uma progressão incremental de seis formulações de fluxo de potência AC: duas formulações de referência baseadas no pacote \textit{PowerModels.jl} (uma de FPO completo e uma de fluxo de potência clássico) e quatro formulações \textit{hand-built} em \textit{JuMP}, nas quais cada uma das ações de controle ANAREDE listadas é introduzida sequencialmente sobre a anterior. As entradas são lidas a partir de arquivos no formato PWF utilizando o pacote \textit{PWF.jl}, com extração das estruturas de controle (DOPC e DLIN) habilitada. O modelo não linear é resolvido com o \textit{solver} Ipopt e os controles são introduzidos via variáveis de folga penalizadas em uma função objetivo do tipo \textit{soft-constraint}. As seis formulações são submetidas a uma bateria padronizada de vinte e sete casos de teste, abrangendo redes pedagógicas de três a nove barras, redes de médio porte (300 e 500 barras), recortes reduzidos do subsistema Nordeste do SIN (\texttt{caso\_red} e \texttt{caso\_red2}) e o \textit{snapshot} integral do Sistema Interligado Nacional (caso \texttt{CASO\_VER\_MAXDIU.PWF} do ciclo PAR/PEL 2027-2031 do ONS, Verão 2027/2028 Máxima Diurna, com 13\,338 barras). Os resultados mostram que as formulações \textit{hand-built} convergem em uma faixa estritamente maior de casos do que as referências \textit{PowerModels}. Em particular, recuperam soluções factíveis em casos com QLIM ativo nos quais o FPO clássico declara o problema infactível, e as soluções produzidas pelas formulações com CSCA e CTAP se agrupam em \textit{clusters} numericamente distintos da referência, reproduzindo o comportamento esperado das heurísticas de controle do ANAREDE.

% Separe as palavras-chave por ponto
\palavraschave{Fluxo de potência ótimo. Sistemas elétricos de potência. JuMP. PowerModels.jl. Ações de controle. Ipopt.}
